\documentclass[11pt]{article}
\usepackage[a4paper,margin=1.9cm]{geometry}
\usepackage{fontspec}
\setmainfont{Latin Modern Roman}
\usepackage{titlesec}
\usepackage{enumitem}
\usepackage[table]{xcolor}
\usepackage{longtable}
\usepackage{booktabs}
\usepackage{array}
\usepackage{pifont}
\usepackage{amssymb}
\usepackage{tikz}
\usepackage{tcolorbox}
\usepackage{hyperref}
\usepackage{fancyhdr}
\usepackage{parskip}

\definecolor{gxteal}{HTML}{0F766E}
\definecolor{gxcyan}{HTML}{0284C7}
\definecolor{gxslate}{HTML}{334155}
\definecolor{gxamber}{HTML}{B45309}
\definecolor{gxred}{HTML}{B91C1C}
\definecolor{gxgreen}{HTML}{15803D}
\definecolor{gxviolet}{HTML}{6D28D9}

\hypersetup{colorlinks=true, linkcolor=gxteal, urlcolor=gxcyan,
  pdftitle={GlobeX -- PS 26132 Model Build Specifications}, pdfauthor={GlobeX Team}}

\titleformat{\section}{\large\bfseries\color{gxteal}}{\thesection}{0.6em}{}
\titleformat{\subsection}{\normalsize\bfseries\color{gxslate}}{\thesubsection}{0.6em}{}
\titlespacing*{\section}{0pt}{1.7em}{0.7em}

\pagestyle{fancy}
\fancyhf{}
\renewcommand{\headrulewidth}{0.4pt}
\fancyhead[L]{\small\color{gxslate}GlobeX --- PS 26132 Model Specs}
\fancyhead[R]{\small\color{gxslate}Built from scratch}
\fancyfoot[C]{\small\thepage}

\newcommand{\yes}{{\color{gxgreen}\ding{108}}}
\newcommand{\pmid}{{\color{gxcyan}\ding{109}}}
\newcommand{\non}{{\color{black!22}\ding{109}}}
\newcommand{\bld}{{\color{gxamber}\ding{53}}}

\newtcolorbox{callout}[1]{colback=gxteal!5, colframe=gxteal!55, boxrule=0.6pt, arc=2pt,
  left=6pt, right=6pt, top=5pt, bottom=5pt, fonttitle=\bfseries\small, coltitle=white,
  colbacktitle=gxteal, title={#1}}
\newtcolorbox{uspbox}[1]{colback=gxviolet!5, colframe=gxviolet!55, boxrule=0.6pt, arc=2pt,
  left=6pt, right=6pt, top=5pt, bottom=5pt, fonttitle=\bfseries\small, coltitle=white,
  colbacktitle=gxviolet, title={#1}}
\newtcolorbox{warnbox}[1]{colback=gxamber!6, colframe=gxamber!65, boxrule=0.6pt, arc=2pt,
  left=6pt, right=6pt, top=5pt, bottom=5pt, fonttitle=\bfseries\small, coltitle=white,
  colbacktitle=gxamber, title={#1}}
\newtcolorbox{redbox}[1]{colback=gxred!6, colframe=gxred!65, boxrule=0.6pt, arc=2pt,
  left=6pt, right=6pt, top=5pt, bottom=5pt, fonttitle=\bfseries\small, coltitle=white,
  colbacktitle=gxred, title={#1}}

\begin{document}

% ================= TITLE =================
\begin{titlepage}
\centering
\vspace*{1.6cm}
{\Huge\bfseries\color{gxteal} PS 26132 --- Model Build}\\[0.4em]
{\Huge\bfseries\color{gxteal} Specifications}\\[1.8em]
{\Large\color{gxslate} What actually needs ML or CV, what needs plain math instead,}\\[0.3em]
{\Large\color{gxslate} and the full build spec for each --- for a team building from scratch}\\[3em]
{\normalsize\color{gxslate} Market linkages \& price discovery for farmers --- Govt.\ of Maharashtra (MSInS)}\\[0.3em]
{\normalsize\color{gxslate} Prepared for GlobeX --- 5 September 2026}\\[2.6em]
\begin{tabular}{ccc}
\textbf{\color{gxcyan}\Large 2} & \textbf{\color{gxgreen}\Large 3} & \textbf{\color{gxamber}\Large 1} \\
{\small components need real ML/CV} & {\small components need only formulas} & {\small stretch upgrade path} \\
\end{tabular}
\vfill
{\footnotesize\color{gxslate} Internal working document --- not for external distribution as-is}
\end{titlepage}

% ================= CLASSIFICATION =================
\section{The classification, at a glance}
Building everything as a trained model is not more rigorous --- it's often less defensible, since a formula
you can explain in one sentence beats a model you can't justify with the data you actually have. Five
capabilities, one honest verdict each.

\renewcommand{\arraystretch}{1.35}
\noindent\begin{longtable}{@{}p{3.6cm}p{2.6cm}p{8.3cm}@{}}
\toprule
\textbf{Capability} & \textbf{Verdict} & \textbf{Why} \\
\midrule
\endhead
Price / sale-window forecast & \textbf{Classical ML, not DL} & A real time-series prediction problem, but
the deployed Dual-Head GRU already lost to a plain moving average on your own benchmark (54.49\% vs.\
29.42\% MAPE). Gradient-boosted quantile regression is the evidence-backed middle tier --- see Section 2. \\[0.4em]
Sale-window ranking (hold/sell, which market) & \textbf{Formula, not ML} & Few, interpretable factors
(price trend, storage cost, transport cost, glut risk); no labelled ``correct decision'' data exists to
train against; a weighted MCDA score is fully explainable. See Section 4. \\[0.4em]
Buyer/FPO matching & \textbf{Formula, for now} & Cold-start problem --- no historical match-outcome
labels exist yet to learn from. Rule-based filtering plus a similarity score outperforms an untrained
learned ranker. Revisit as ML once real outcomes accumulate. See Section 4. \\[0.4em]
Trust/reputation score & \textbf{Formula, on purpose} & Must be independently recomputable by anyone from
the trade history --- that auditability is the entire point, and a trained model's score can't be verified
the same way. See Section 4. \\[0.4em]
Quality grading & \textbf{Classical CV + simple classifier; DL is a funded stretch} & Grade genuinely
depends on visual attributes text can't capture, so this is the one place CV earns its place --- but there
is no public labelled dataset of AGMARK-graded produce photos, so treat this as a data-collection problem
first. See Section 3. \\
\bottomrule
\end{longtable}

% ================= FORECASTING =================
\clearpage
\section{Spec: price / sale-window forecaster}

\subsection*{Why classical ML, not deep learning}
This isn't a guess --- it matches the published literature on agricultural price forecasting. Recent work
comparing tree-based and hybrid models against pure deep-learning approaches (e.g.\ a TCN--XGBoost hybrid
for agricultural market prices, and broader comparisons of linear/tree-based ML against commodity price
series) consistently finds gradient-boosted trees competitive with or better than deep sequence models on
this kind of data, especially when the historical series per commodity--market pair is not very long ---
exactly your situation with 2--3 states of mandi history. Deep learning needs volume classical ML doesn't.

\subsection*{Build spec}
\renewcommand{\arraystretch}{1.3}
\noindent\begin{longtable}{@{}p{3.0cm}p{9.5cm}@{}}
\toprule
\textbf{Element} & \textbf{Specification} \\
\midrule
\endhead
Target variable & Next-$N$-day modal price per (commodity, market) pair, plus the 10th/50th/90th percentile
bands (quantile targets), not a single point forecast. \\[0.3em]
Features & Lagged prices (1/7/14/30-day), rolling mean/std of price and arrivals, day-of-week and
day-of-year seasonality encodings, arrival-volume trend, min--max price spread (volatility signal ---
unique to mandi data, keep it), state/market one-hot or target-encoded identifiers. \\[0.3em]
Algorithm & Gradient-boosted quantile regression (XGBoost/LightGBM with quantile or pinball-loss
objective), one model per quantile (10th/50th/90th) or a native multi-quantile objective if the library
supports it. \\[0.3em]
Baselines to beat & 7-day and 30-day moving average; naive last-value carry-forward; simple linear trend.
Report all three alongside the model --- this is exactly the comparison your own benchmark file already
does, keep doing it honestly. \\[0.3em]
Train/test split & Chronological, never random --- train on earlier dates, test on later ones, per
commodity--market pair, mirroring the anomaly classifier's existing chronological-split discipline. \\[0.3em]
Evaluation metric & Pinball loss (quantile loss) for the bands; MAPE and R\textsuperscript{2} on the median
prediction for headline comparability against the baselines. \\[0.3em]
What ``good'' looks like & Beats the moving-average baseline on MAPE \emph{and} the 80\% interval (10th--90th
percentile band) actually contains the true price roughly 80\% of the time on held-out data --- calibration
matters as much as point accuracy here, since the product promise is a confidence band, not a number. \\[0.3em]
Stretch upgrade & A hybrid architecture (tree-based model plus a lightweight sequence model for residuals)
only if the median-model baseline is solidly beaten first and time remains --- don't reach for this before
the simple version is honestly evaluated. \\
\bottomrule
\end{longtable}

% ================= QUALITY GRADING =================
\clearpage
\section{Spec: quality grading (computer vision)}

\subsection*{The honest starting point: no ready-made labelled dataset exists}
Public fruit/vegetable image datasets do exist on Kaggle (fresh-vs-rotten classification, general
fruit-type recognition, some defect/quality labels), and there is a real, peer-reviewed literature on
computer-vision produce grading using exactly the classical features below. But none of the public datasets
are labelled against AGMARK's actual grade categories (moisture, size, foreign-matter thresholds per
commodity) --- they're closest as a \emph{pretraining or sanity-check} resource, not a drop-in substitute.
Budget real time for photographing and labelling your own small sample set; that is the actual bottleneck,
not the modelling.

\subsection*{Tiered build plan --- start at Tier 1, not Tier 3}
\renewcommand{\arraystretch}{1.35}
\noindent\begin{longtable}{@{}p{1.6cm}p{4.4cm}p{6.5cm}@{}}
\toprule
\textbf{Tier} & \textbf{Approach} & \textbf{Why this order} \\
\midrule
\endhead
0 (baseline, always keep) & OCR/NER check of the submitted grading certificate against the lot record &
Already specified, already reliable, zero training data needed --- this is the guaranteed-working fallback
if the CV tiers below run out of time. \\[0.4em]
1 (build this) & Classical image features --- size via a reference object in frame, colour histogram
compared against AGMARK's published per-commodity thresholds, blob/contour-based defect and blemish
counting via OpenCV --- feeding a simple classifier (Random Forest or SVM) & Data-efficient: works with as
few as 50--150 hand-labelled photos per commodity, fully explainable (``flagged for X\% surface blemish
area,'' not a black-box score), and matches the established computer-vision-for-produce-grading literature
rather than inventing an untested approach. \\[0.4em]
2 (stretch only) & Small transfer-learning CNN (e.g.\ fine-tuned MobileNet/EfficientNet) trained on your
own labelled photos & Only attempt if Tier 1 is working and you can realistically gather several hundred
labelled images per class --- a CNN trained on too few images will overfit and perform \emph{worse} than
Tier 1, which is a real risk if attempted under time pressure. \\
\bottomrule
\end{longtable}

\subsection*{Build spec (Tier 1)}
\renewcommand{\arraystretch}{1.3}
\noindent\begin{longtable}{@{}p{3.0cm}p{9.5cm}@{}}
\toprule
\textbf{Element} & \textbf{Specification} \\
\midrule
\endhead
Target variable & Grade category per commodity (e.g.\ AGMARK Grade A/B/C, or a pass/fail against the
declared certificate grade), a multi-class or ordinal classification target. \\[0.3em]
Features & Estimated size/diameter (calibrated against a reference object of known size in the photo),
colour histogram distance from the expected healthy-produce range, percentage of surface area flagged as
blemished/discoloured via contour/blob detection, count of visible defects. \\[0.3em]
Algorithm & Random Forest or SVM on the engineered features above --- deliberately not a CNN at this tier. \\[0.3em]
Data plan & Self-photograph 50--150 samples per commodity across the grade categories you intend to
support (start with 1--2 commodities, not all of them); optionally pretrain feature-extraction thresholds
using a public freshness/defect Kaggle dataset as a sanity check, but do not treat it as ground truth for
AGMARK grades. \\[0.3em]
Train/test split & Stratified by grade class, held-out test images never touched during threshold-tuning. \\[0.3em]
Evaluation metric & Per-class precision/recall/F1, plus a confusion matrix --- report where it actually
fails (which grade pairs get confused), don't just report overall accuracy. \\[0.3em]
What ``good'' looks like & Clearly beats a trivial baseline (always-predict-the-most-common-grade); more
importantly, its mistakes should be between \emph{adjacent} grades (B misclassified as A or C), not wild
misses --- that's the difference between ``imperfect but usable'' and ``untrustworthy.'' \\
\bottomrule
\end{longtable}

% ================= FORMULAS =================
\clearpage
\section{Spec: the three formula-based components}
These are specified, not trained --- validate them like you would any deterministic system: with test
cases and backtests, not train/test splits.

\subsection*{Sale-window ranking (MCDA)}
\textbf{Inputs:} price-trend score, storage-cost estimate, transport-cost estimate, arrival-volume/glut-risk
score --- one row per candidate action (sell now at market $X$, hold $N$ days). \textbf{Logic:} weighted sum
(or TOPSIS-style normalised distance) across the four dimensions; weights set from domain judgement
initially, refined by backtesting against historical price outcomes once enough mandi history is ingested
(the same backtesting discipline the existing ranking engine already used to reach $\rho = 0.884$).
\textbf{Output:} a ranked list of actions with the top choice and its reasoning (which dimension drove the
recommendation) shown, not just a number. \textbf{Validation:} backtest --- for past dates where the actual
future price is now known, check whether the engine's top-ranked action would have been the better one more
often than a naive ``always sell immediately'' baseline.

\subsection*{Buyer/FPO matching}
\textbf{Inputs:} commodity, grade tier, volume range, location, (once available) buyer's on-chain trust
score. \textbf{Logic:} hard filters first (same commodity, compatible grade tier, volume within range,
location within radius), then a weighted similarity score across the remaining soft factors, ranked
descending. \textbf{Output:} a ranked shortlist of candidate counterparties, not a single forced match.
\textbf{Validation:} manual spot-checks against the pilot seed dataset now; once real trades happen,
graduate to precision@$k$ against actual accepted matches.

\subsection*{Trust/reputation score}
\textbf{Inputs:} completed, disputed, cancelled, delivered, inspected, on-time, quality-passed trade counts
--- exactly the fields \texttt{TradeLedger.sol} already aggregates on-chain. \textbf{Logic:} a transparent
weighted formula (e.g.\ completion rate, on-time rate and quality-pass rate combined, disputes penalised) that
anyone can recompute by hand from the public trade history. \textbf{Output:} a single score plus the
component rates that produced it, always shown together --- never the score alone. \textbf{Validation:}
the formula must be independently reproducible from on-chain data by a third party; that reproducibility
\emph{is} the test, not a held-out accuracy number.

% ================= SUMMARY =================
\clearpage
\section{Build order}
\begin{enumerate}[leftmargin=1.6em]
\item Formula-based components first (ranking, matching, trust score) --- fastest to build, zero training
data required, and they unblock the rest of the product flow immediately.
\item Forecaster (Tier: classical ML) --- needs real ingested mandi data first (see the dataset landscape
doc), then a day or two of feature engineering and baseline comparison.
\item Quality grading Tier 0 (OCR/certificate check) --- reuse existing document-intelligence patterns,
low effort.
\item Quality grading Tier 1 (classical CV + simple classifier) --- start photographing and labelling
samples as early as possible; this is the genuine bottleneck in the whole plan, not the modelling step.
\item Quality grading Tier 2 (CNN) --- attempt only if Tier 1 works and time plus labelled images allow.
\end{enumerate}

\begin{warnbox}{The point of this document}
Two components get real modelling effort because the problem genuinely needs statistical learning
(forecasting) or visual pattern recognition (grading). Three components deliberately don't get a trained
model, because a transparent formula is both more honest about what data you actually have and more
defensible in front of a judge who asks ``how does this work'' and expects a real answer, not ``the model
learned it.''
\end{warnbox}

\end{document}
